% =============================================================================
%  Chapter 60 — Steppable Dynamic Analysis with StepDirector
% =============================================================================

\chapter{Steppable Dynamic Analysis with StepDirector}
\label{ch:steppable-dynamic}

This chapter documents the design, implementation, and testing of the
single-step API for \texttt{DynamicAnalysis} and the \texttt{StepDirector}
abstraction for condition-based breakpoints.  Together these enable
pausing a dynamic simulation at specified conditions, inspecting or
modifying the model, changing integration parameters, and resuming ---
all without modifying the analysis class itself.

% =============================================================================
\section{Motivation}
\label{sec:sd-motivation}

The batch \texttt{solve(t\_final, dt)} interface of \texttt{DynamicAnalysis}
delegates the entire time-stepping loop to PETSc's \texttt{TSSolve}, which
runs from $t=0$ to $t_\text{final}$ in a single call.  While convenient,
this monolithic interface prevents several important use cases:

\begin{itemize}
  \item \textbf{Pause \& inspect:}  Capture the solution at specific times
        to couple with external models, extract quantities of interest,
        or generate intermediate output.
  \item \textbf{Runtime reconfiguration:}  Change the time step~$\Delta t$,
        switch the integration scheme (e.g.\ from implicit
        Generalized-$\alpha$ to explicit Euler), or modify boundary
        conditions mid-simulation.
  \item \textbf{Multi-phase analysis:}  Run a static preload phase
        (via \texttt{NonlinearAnalysis}), transfer the state, and
        continue with a dynamic phase with different control logic.
  \item \textbf{Condition-based stopping:}  Terminate or pause when
        a displacement exceeds a threshold, a damage variable reaches
        a limit, or an external convergence criterion is met.
\end{itemize}

The design goal is to provide these capabilities with zero overhead on the
existing batch API and minimal modifications to the class.

% =============================================================================
\section{Design}
\label{sec:sd-design}

\subsection{Single-Step API on DynamicAnalysis}

Three new public methods expose single-step control:

\begin{lstlisting}
bool step();                                // one time step
StepVerdict step_to(double t_target);       // advance to t
StepVerdict step_n(int n);                  // advance n steps
\end{lstlisting}

\noindent
\texttt{step()} wraps PETSc's \texttt{TSStep(ts\_)} with the bookkeeping
required for manual single-step control:

\begin{enumerate}
  \item Reset the convergence reason to \verb|TS_CONVERGED_ITERATING|
        (PETSc refuses to step if the TS is flagged as already converged).
  \item Set the step budget to \verb|current_step() + 1| via
        \verb|TSSetMaxSteps|.
  \item Set a sufficiently large \verb|TSSetMaxTime| to avoid
        time-limit convergence.
  \item Call \verb|TSStep(ts_)|.
  \item Check \verb|TSGetConvergedReason|: return \texttt{true} if
        the reason is non-negative (converged or max-steps reached),
        \texttt{false} on divergence.
\end{enumerate}

The Monitor callback registered in \texttt{setup()} is invoked
automatically by PETSc during \texttt{TSStep}, so the full post-step
pipeline (material commit, model state update, observer notification)
runs without duplication.

\subsection{Extracted \texttt{post\_step\_} Method}

The material commit, model state vector update, and observer/monitor
notification logic was extracted from the static \texttt{Monitor} callback
into a private member function:

\begin{lstlisting}
void post_step_(PetscInt step, double t, Vec U, Vec V);
\end{lstlisting}

\noindent
The static \texttt{Monitor} callback becomes a thin trampoline:

\begin{lstlisting}
static PetscErrorCode Monitor(TS ts, PetscInt step,
                              PetscReal t, Vec U, void* ctx)
{
    auto* self = static_cast<Context*>(ctx)->self;
    Vec V;
    TS2GetSolution(ts, &U, &V);
    self->post_step_(step, t, U, V);
    return PETSC_SUCCESS;
}
\end{lstlisting}

This refactoring keeps the Monitor API surface unchanged while enabling
direct invocation from future code paths where the PETSc monitor is not
active.

\subsection{Runtime Reconfiguration}

Four additional methods allow runtime parameter changes between steps:

\begin{lstlisting}
void   set_time_step(double dt);
double get_time_step() const;
void   set_integration_method(TSType type);
void   set_director(StepDirector<ModelT> dir);
\end{lstlisting}

\noindent
\texttt{set\_integration\_method} wraps \texttt{TSSetType} and allows
switching between \texttt{TSALPHA2}, \texttt{TSNEWMARK}, \texttt{TSEULER},
\texttt{TSRK}, etc.\ at runtime.

% =============================================================================
\section{StepDirector — Condition-Based Breakpoints}
\label{sec:step-director}

\subsection{StepVerdict Enum}

The director protocol uses a three-valued verdict:

\begin{lstlisting}
enum class StepVerdict {
    Continue,   // keep stepping
    Pause,      // return control to caller (resumable)
    Stop        // terminate the analysis
};
\end{lstlisting}

\noindent
The ordering $\text{Continue} < \text{Pause} < \text{Stop}$ enables a
\texttt{most\_restrictive(a, b)} combinator for composing multiple
directors.

\subsection{StepDirector Type Alias}

A \texttt{StepDirector<ModelT>} is a type-erased callable:

\begin{lstlisting}
template <typename ModelT>
using StepDirector = std::function<
    StepVerdict(const StepEvent&, const ModelT&)>;
\end{lstlisting}

\noindent
The director is evaluated \emph{after} every converged step in
\texttt{step\_to()} and \texttt{step\_n()}.  This guarantees the
model is in a fully consistent, committed state when the caller
inspects it upon a \texttt{Pause} verdict.

\subsection{Convenience Factories}

The \texttt{step\_director} namespace provides ready-made directors:

\begin{center}
\begin{tabular}{ll}
  \toprule
  Factory & Description \\
  \midrule
  \texttt{pause\_at\_times(\{$t_1, t_2, \ldots$\})} &
    Pause when $t \geq t_i$ (each breakpoint fires once) \\
  \texttt{pause\_every\_n($N$)} &
    Pause every $N$ steps \\
  \texttt{pause\_on($F$)} &
    Pause when predicate $F(\text{ev}, \text{model})$ is true \\
  \texttt{stop\_on($F$)} &
    Stop when predicate $F(\text{ev}, \text{model})$ is true \\
  \texttt{compose($d_1, d_2, \ldots$)} &
    Most restrictive verdict of all directors \\
  \bottomrule
\end{tabular}
\end{center}

\subsection{Integration with DynamicAnalysis}

The \texttt{step\_to} and \texttt{step\_n} methods accept an optional
director parameter.  A persistent director can also be set via
\texttt{set\_director()}.

\begin{lstlisting}
// Pause at specific times
auto dir = step_director::pause_at_times<MyModel>({1.0, 2.5});
auto verdict = solver.step_to(10.0, dir);
if (verdict == StepVerdict::Pause) {
    // inspect model, modify BCs, change dt, etc.
    solver.step_to(10.0, dir);   // resume
}
\end{lstlisting}

% =============================================================================
\section{Backward Compatibility}
\label{sec:sd-backward}

The existing \texttt{solve(t\_final, dt)} and
\texttt{solve(t\_start, t\_final, dt)} methods remain unchanged and
continue to use PETSc's \texttt{TSSolve} internally.  No existing
call site is affected.  The Monitor callback extraction is a pure
refactoring with no behavioral change for the batch API.

Test~9 in the verification suite confirms that \texttt{step\_to(t)}
produces results \emph{identical} to \texttt{solve(t, dt)} (difference
$= 0$), validating that the single-step pathway is numerically
equivalent to the batch pathway.

% =============================================================================
\section{Verification}
\label{sec:sd-verification}

A dedicated test suite (\texttt{test\_steppable\_dynamic.cpp}) validates
10 scenarios with 32 assertions:

\begin{center}
\begin{tabular}{clc}
  \toprule
  Test & Description & Assertions \\
  \midrule
  1  & \texttt{step()} — single time step convergence & 5 \\
  2  & \texttt{step\_n(10)} — exact step count         & 3 \\
  3  & \texttt{step\_to(t)} — target time reached       & 3 \\
  4  & \texttt{pause\_at\_times} breakpoint             & 5 \\
  5  & \texttt{pause\_every\_n} periodic pause          & 4 \\
  6  & \texttt{pause\_on} custom condition              & 2 \\
  7  & \texttt{compose} directors                       & 2 \\
  8  & Runtime $\Delta t$ change                        & 3 \\
  9  & Equivalence: step-by-step = batch solve          & 3 \\
  10 & \texttt{stop\_on} termination                    & 2 \\
  \midrule
     & \textbf{Total}                                   & \textbf{32} \\
  \bottomrule
\end{tabular}
\end{center}

All 32 assertions pass.  Additionally, the full existing test suite
(46/50, with 4 pre-existing failures due to missing data files) passes
unchanged, confirming backward compatibility.

% =============================================================================
\section{Files Modified and Created}
\label{sec:sd-files}

\begin{center}
\begin{tabular}{lll}
  \toprule
  File & Action & Purpose \\
  \midrule
  \texttt{src/analysis/StepDirector.hh} & Created &
    \texttt{StepVerdict} + \texttt{StepDirector<ModelT>} + factories \\
  \texttt{src/analysis/DynamicAnalysis.hh} & Modified &
    \texttt{post\_step\_()}, \texttt{step()}, \texttt{step\_to()},
    \texttt{step\_n()}, reconfig \\
  \texttt{src/analysis/Analysis.hh} & Modified &
    Added \texttt{StepDirector.hh} include \\
  \texttt{tests/test\_steppable\_dynamic.cpp} & Created &
    10 tests, 32 assertions \\
  \texttt{CMakeLists.txt} & Modified &
    Registered \texttt{fall\_n\_steppable\_dynamic\_test} \\
  \bottomrule
\end{tabular}
\end{center}

% =============================================================================
\section{Future Directions}
\label{sec:sd-future}

This infrastructure enables the remaining phases of the unified
steppable solver plan:

\begin{description}
  \item[Phase C — SteppableSolver concept:]  Unify
    \texttt{NonlinearAnalysis} and \texttt{DynamicAnalysis} under a
    common concept with \texttt{step()}, \texttt{step\_to()}, and
    \texttt{current\_time()}.  An \texttt{AnalysisState} struct will
    enable cross-engine state transfer.
  \item[Phase D — AnalysisDirector:]  A high-level orchestrator that
    sequences multiple analysis phases (static preload $\to$ dynamic
    excitation $\to$ static gravity redistribution), with automatic
    state transfer and observer lifecycle management.
\end{description}
